\chapter{Anéis e aritmética}\label{ch:b3:rings}

Os inteiros usuais se fatoram de maneira única em primos; o mesmo
ocorre com os polinômios sobre um corpo. Esses dois fatos são um único
teorema? Este capítulo responde que sim e encontra as hipóteses exatas
que tornam possível uma ``aritmética'' em um anel comutativo: a cadeia
\[
\text{euclidiana} \;\Longrightarrow\; \text{principal}
\;\Longrightarrow\; \text{fatorial (DFU)},
\]
com todas as implicações demonstradas e todas as recíprocas refutadas.
A teoria é então posta à prova onde ela mostra seu valor: os inteiros
de Gauss $\Z[\iu]$ (que abrirão o teorema dos dois quadrados de Fermat no
problema de fim de semana), os anéis de polinômios em várias variáveis
(lema de Gauss, critério de \omterm{thm:b3:rings:criteria}{Eisenstein}) e os anéis \omterm{def:b3:rings:noetherian}{noetherianos},
culminando no teorema da base de Hilbert. Ao longo de todo o capítulo,
\emph{anel} significa anel comutativo com unidade $1 \neq 0$; os \omterm{def:b3:rings:ideal}{ideais} de
$\Z$ e de $K[X]$ vistos no volume do segundo ano são nossos dois
exemplos condutores.

\section{Ideais, quocientes e o teorema de isomorfismo}

\begin{definition}\label{def:b3:rings:ideal}
Um \emph{ideal}\index{ideal} $I$ de um anel $A$ é um subgrupo
aditivo tal que $AI \subseteq I$. O \emph{anel quociente}\index{anel quociente}
$A/I$ é o \omterm{thm:b3:groups:quotient}{grupo quociente} $(A, +)/I$ munido da multiplicação $(a + I)(b + I) = ab + I$: ela está
bem definida, pois trocar $a$ por $a + x$ e $b$ por
$b + y$ ($x, y \in I$) altera $ab$ de $ay + xb + xy \in I$. A projeção $\pi \colon A
\to A/I$ é um morfismo
de anéis sobrejetor de núcleo $I$, e os núcleos dos morfismos de
anéis são exatamente os ideais.
\end{definition}

\begin{theorem}[Primeiro teorema de isomorfismo]\label{thm:b3:rings:firstiso}
Se $f \colon A \to B$ é um morfismo de anéis, então $\bar f\colon
A/\ker f \to \operatorname{im} f$, $a + \ker f \mapsto f(a)$, é um isomorfismo de
anéis. Mais geralmente, $f$ se fatora por $A/I$ para todo \omterm{def:b3:rings:ideal}{ideal}
$I \subseteq \ker f$. Os \omterm{def:b3:rings:ideal}{ideais} de $A/I$ são os $J/I$ com $J \supseteq I$ \omterm{def:b3:rings:ideal}{ideal} de $A$
(teorema da correspondência).
\end{theorem}

\begin{proof}
Como no caso dos grupos (\cref{thm:b3:groups:firstiso,%
thm:b3:groups:correspondence}), observando que todas as aplicações em
jogo respeitam também os produtos: $\bar f$ está bem definida, é
bijetora sobre a imagem e multiplicativa; a correspondência $J \mapsto J/I$, $\bar J \mapsto \pi^{-1}(\bar J)$
preserva \omterm{def:b3:rings:ideal}{ideais} nos dois sentidos porque $\pi$ é um morfismo de
anéis sobrejetor.
\end{proof}

\begin{definition}\label{def:b3:rings:primemaximal}
Seja $I \subsetneq A$ um \omterm{def:b3:rings:ideal}{ideal} próprio. $I$ é \emph{primo}\index{ideal primo}
se $ab \in I \Rightarrow a \in I$ ou $b \in I$; $I$ é \emph{maximal}\index{ideal maximal} se
nenhum \omterm{def:b3:rings:ideal}{ideal} está estritamente entre $I$ e $A$.
\end{definition}

\begin{proposition}\label{prop:b3:rings:primemaximal}
$I$ é primo $\iff$ $A/I$ é um domínio de integridade; $I$ é
\omterm{def:b3:rings:primemaximal}{maximal} $\iff$ $A/I$ é um corpo. Em particular, os \omterm{def:b3:rings:primemaximal}{ideais
maximais} são primos.
\end{proposition}

\begin{proof}
Escreva $\bar a$ para as classes em $A/I$. ``$I$ primo''
traduz-se literalmente por ``$\bar a\bar b = 0 \Rightarrow \bar a = 0$ ou $\bar b
= 0$'', e $A/I \neq 0$ por $I \neq A$: essa é a
definição de domínio. Para a maximalidade, use o teorema da
correspondência: nenhum \omterm{def:b3:rings:ideal}{ideal} estritamente entre $I$ e $A$
$\iff$ $A/I$ não tem \omterm{def:b3:rings:ideal}{ideal} além de $0$ e de si mesmo
$\iff$ $A/I$ é um corpo --- quanto a esta última etapa: em um
corpo, os únicos \omterm{def:b3:rings:ideal}{ideais} são $0$ e o corpo todo (um \omterm{def:b3:rings:ideal}{ideal} que contém
$x \ne 0$ contém $x^{-1}x = 1$); reciprocamente, se todo $x$ não nulo gera o \omterm{def:b3:rings:ideal}{ideal}
unidade, então $xy = 1$ para algum $y$. Corpos são domínios, logo
os \omterm{def:b3:rings:primemaximal}{ideais maximais} são primos.
\end{proof}

\begin{example}\label{ex:b3:rings:primemaximal}
Em $\Z$: os \omterm{def:b3:rings:primemaximal}{ideais primos} são $(0)$ e os $(p)$, com $p$
primo; os maximais são os $(p)$ ($\Z/p\Z = \mathbb F_p$ é um corpo, $\Z/(0) = \Z$ não é). Em
$K[X, Y]$: $(X) \subsetneq (X, Y)$ são ambos primos ($K[X,Y]/(X) \cong K[Y]$, um domínio; $K[X,Y]/(X,Y) \cong K$, um corpo), de modo
que $(X)$ é primo mas não \omterm{def:b3:rings:primemaximal}{maximal}.
\end{example}

Para garantir em plena generalidade que \omterm{def:b3:rings:primemaximal}{ideais maximais} \emph{existem},
precisamos de um princípio conjuntista. Um conjunto parcialmente
ordenado é \emph{indutivo} se todo subconjunto totalmente ordenado
(\emph{cadeia}) admite uma cota superior.

\begin{theorem}[Lema de Zorn]\label{thm:b3:rings:zorn}
Todo conjunto parcialmente ordenado indutivo e não vazio tem um
elemento \omterm{def:b3:rings:primemaximal}{maximal}.\index{lema de Zorn}
\admitted
\end{theorem}

\begin{remark}\label{rem:b3:rings:zorn}
Isso não é um teorema da matemática usual, mas um \emph{axioma}: ele é
equivalente, sobre os axiomas básicos de Zermelo--Fraenkel da teoria dos
conjuntos, ao axioma da escolha (``todo produto de conjuntos não vazios
é não vazio''), que aceitamos ao longo de todo este livro. Sinalizamos
cada uso. A análise voltará a invocá-lo (Hahn--Banach,
\cref{ch:b3:banach}).
\end{remark}

\begin{theorem}[Krull]\label{thm:b3:rings:krull}
Todo \omterm{def:b3:rings:ideal}{ideal} próprio $I \subsetneq A$ está contido em um \omterm{def:b3:rings:primemaximal}{ideal maximal}.\index{teorema
de Krull}
\end{theorem}

\begin{proof}
Ordene por inclusão o conjunto $\mathcal E$ dos \omterm{def:b3:rings:ideal}{ideais} próprios que contêm
$I$; ele é não vazio ($I \in \mathcal E$). Uma cadeia $(J_\lambda)$ em $\mathcal E$ tem por cota
superior $J = \bigcup
J_\lambda$: um \omterm{def:b3:rings:ideal}{ideal} (dois elementos $a, b \in J$ quaisquer estão em um
mesmo $J_\lambda$, por totalidade), próprio ($1 \notin J_\lambda$ para todo $\lambda$) e que
contém $I$. O lema de Zorn fornece um elemento \omterm{def:b3:rings:primemaximal}{maximal} de $\mathcal
E$,
que é um \omterm{def:b3:rings:primemaximal}{ideal maximal} contendo $I$ (um \omterm{def:b3:rings:ideal}{ideal} próprio estritamente
acima dele estaria em $\mathcal E$).
\end{proof}

\begin{theorem}[Teorema chinês dos restos]\label{thm:b3:rings:crt}
Sejam $I_1, \dots, I_n$ \omterm{def:b3:rings:ideal}{ideais} de $A$ dois a dois
\emph{comaximais}\index{ideais comaximais} ($I_k + I_l = A$ para $k \neq l$). Então
\[
A\Big/\bigcap_{k=1}^n I_k \;\xrightarrow{\;\sim\;}\;
\prod_{k=1}^n A/I_k,
\qquad
a \longmapsto (a + I_1, \dots, a + I_n),
\]
e, além disso, $\bigcap_k I_k = I_1 I_2 \cdots I_n$ (o \omterm{def:b3:rings:ideal}{ideal} gerado pelos produtos).\index{teorema
chinês dos restos}
\end{theorem}

\begin{proof}
A aplicação $f(a) = (a + I_k)_k$ é um morfismo de anéis de núcleo $\bigcap I_k$; pelo~\cref{thm:b3:rings:firstiso},
basta demonstrar a sobrejetividade. Fixe $k$; para cada $l \neq k$
escreva $1 = u_l +
v_l$ com $u_l \in I_k$, $v_l \in I_l$ (comaximalidade). Então
\[
e_k = \prod_{l \neq k} v_l = \prod_{l\neq k}(1 - u_l)
\equiv 1 \pmod{I_k},
\qquad e_k \in I_l \ (l \neq k),
\]
logo $f(e_k) = (0, \dots, 1, \dots, 0)$; dado um alvo $(a_k +
I_k)_k$, o elemento $\sum_k a_k e_k$ é levado nele.

Produtos e interseção: $I_1\cdots I_n \subseteq \bigcap I_k$ sempre. Reciprocamente, por indução basta
tratar o caso $n = 2$ (verifica-se que $I_1$ e $I_2\cdots I_n$ são
\omterm{thm:b3:rings:crt}{comaximais}: multiplicar $1 = u_l + v_l$ para $l \geq 2$ dá $1 \in I_1 + I_2\cdots
I_n$). Para $n = 2$:
escreva $1 = u + v$, $u \in I_1$, $v \in I_2$; para $x \in I_1 \cap I_2$, $x = xu + xv \in I_2I_1 + I_1I_2 =
I_1I_2$.
\end{proof}

\begin{example}\label{ex:b3:rings:crtz}
Em $\Z$ com $I_k = (m_k)$ e $m_k$ dois a dois coprimos: $\Z/(m_1\cdots
m_n)\Z \cong \prod \Z/m_k\Z$ --- o
teorema chinês dos restos do volume do segundo ano. Restringindo às
unidades: $(\Z/mn\Z)^\times \cong
(\Z/m\Z)^\times \times (\Z/n\Z)^\times$ para $\gcd(m,n)=1$, donde a multiplicatividade da função $\varphi$
de Euler (\cref{exo:b3:rings:8}).
\end{example}

\section{Divisibilidade: euclidianos, principais, fatoriais}

\begin{definition}\label{def:b3:rings:divisibility}
Sejam $A$ um domínio de integridade e $a, b \in A$. Dizemos que $a$
\emph{divide} $b$ ($a \mid b$) se $b \in (a) = aA$. Os elementos $a,
b$ são
\emph{associados}\index{elementos associados} se $a = ub$ com $u \in A^\times$
(equivalentemente, $(a) = (b)$). Um elemento $p$ não nulo e não inversível
é:
\begin{itemize}
  \item \emph{irredutível}\index{elemento irredutível} se $p = ab$
        força $a \in A^\times$ ou $b \in A^\times$;
  \item \emph{primo}\index{elemento primo} se $p \mid ab$ força $p
        \mid a$ ou
        $p \mid b$ (isto é, o \omterm{def:b3:rings:ideal}{ideal} $(p)$ é primo).
\end{itemize}
\end{definition}

\begin{proposition}\label{prop:b3:rings:primeirred}
Em todo domínio, primo $\Rightarrow$ \omterm{def:b3:rings:divisibility}{irredutível}. A recíproca é falsa em geral:
em $\Z[\iu\sqrt 5] = \{a + \iu b\sqrt5 : a, b
\in \Z\}$, o elemento $2$ é \omterm{def:b3:rings:divisibility}{irredutível} mas não é primo.
\end{proposition}

\begin{proof}
Seja $p$ primo e $p = ab$. Então $p \mid ab$, digamos $p \mid a$:
$a = pc$, logo $p = pcb$, e cancelando $p$ (é um domínio!) obtém-se
$cb =
1$: $b \in A^\times$.

Em $\Z[\iu\sqrt5]$, use a norma $N(x + \iu y\sqrt 5) = x^2 +
5y^2$, que é multiplicativa (ela vale $\abs z^2$). Se
$2 = ab$ com $a, b$ não inversíveis, então $4 = N(a)N(b)$ com $N(a), N(b) \neq
1$ (os
elementos de norma $1$ são $\pm1$, as unidades), logo $N(a) = 2$:
impossível, pois $x^2 + 5y^2 = 2$ não tem solução inteira. Assim $2$ é
\omterm{def:b3:rings:divisibility}{irredutível}. Mas $2 \mid 6 = (1 + \iu\sqrt5)(1 - \iu\sqrt5)$ enquanto $2$ não divide nenhum dos fatores
($\frac12 \pm \frac{\iu\sqrt5}2
\notin \Z[\iu\sqrt5]$): ele não é primo.
\end{proof}

\begin{definition}\label{def:b3:rings:pidufd}
Um domínio de integridade $A$ é:
\begin{itemize}
  \item \emph{euclidiano}\index{domínio euclidiano} se existe uma
        aplicação $\nu \colon A \setminus \{0\} \to \N$ (uma \emph{função euclidiana}) tal que, para
        todos $a, b$ com $b \ne 0$, existem $q, r$ com $a = bq + r$ e
        ($r = 0$ ou $\nu(r) <
        \nu(b)$);
  \item \emph{principal} (um \emph{DIP}\index{domínio de ideais
        principais}) se todo \omterm{def:b3:rings:ideal}{ideal} é da forma $(a)$;
  \item \emph{fatorial} (um \emph{DFU}\index{domínio de fatoração
        única}) se todo elemento não nulo e não inversível é um produto
        de irredutíveis, único a menos da ordem e de \omterm{def:b3:rings:divisibility}{associados}.
\end{itemize}
\end{definition}

\begin{theorem}\label{thm:b3:rings:euclideanpid}
\omterm{def:b3:rings:pidufd}{Euclidiano} $\Rightarrow$ principal.
\end{theorem}

\begin{proof}
Sejam $I \neq (0)$ um \omterm{def:b3:rings:ideal}{ideal} e $b \in I \setminus\{0\}$ com $\nu(b)$ mínimo. Para $a \in I$,
divida: $a = bq + r$; então $r = a
- bq \in I$, e $\nu(r) < \nu(b)$ contradiria a minimalidade, de modo
que $r = 0$ e $a \in (b)$: $I = (b)$.
\end{proof}

\begin{example}\label{ex:b3:rings:euclidean}
$\Z$ (com $\nu = \abs\cdot$) e $K[X]$ (com $\nu = \deg$) são \omterm{def:b3:rings:pidufd}{euclidianos} --- as
duas divisões foram demonstradas no volume do segundo ano. O mesmo vale
para $\Z[\iu]$, com $\nu = N$ a norma quadrática (\cref{exo:b3:rings:4}); a geometria
da demonstração está na figura abaixo. Existe um \omterm{def:b3:rings:pidufd}{DIP} que não é
\omterm{def:b3:rings:pidufd}{euclidiano}, mas é delicado certificá-lo (o exemplo padrão é $\Z\bigl[\frac{1+\iu\sqrt{19}}2\bigr]$); um
\omterm{def:b3:rings:pidufd}{DFU} que não é \omterm{def:b3:rings:pidufd}{DIP} é fácil: $K[X, Y]$ (\cref{exo:b3:rings:6}), ou $\Z[X]$.
\end{example}

\begin{omfigure}
\begin{tikzpicture}[scale=1.05]
  \clip (-2.7,-1.7) rectangle (2.7,2.2);
  \foreach \x in {-2,...,2}
    \foreach \y in {-1,...,2}
      \fill[black!55] (\x,\y) circle (1.4pt);
  \draw[omDef, thick] (0.6,0.4) circle (1);
  \fill[omThm] (0.6,0.4) circle (1.8pt)
    node[below right, font=\small, omThm] {$a/b$};
  \fill[omDef] (1,0) circle (2pt)
    node[below, font=\small, omDef] {$q$};
  \draw[omDef, ->, shorten >=2.5pt, shorten <=2.5pt]
    (0.6,0.4) -- (1,0);
\end{tikzpicture}

{\small Divisão nos inteiros de Gauss: o quociente exato $a/b \in \C$ está a
distância no máximo $\leq \frac{\sqrt2}{2} < 1$ de algum ponto do reticulado $q \in \Z[\iu]$; então
$r = a - bq$ satisfaz $N(r) =
N(b)\,\abs{a/b - q}^2 < N(b)$. Uma divisão euclidiana e, com ela, toda uma
aritmética.}
\end{omfigure}

\begin{lemma}[Cadeias ascendentes de ideais principais]
\label{lem:b3:rings:acc}
Em um \omterm{def:b3:rings:pidufd}{DIP}, toda sequência crescente de \omterm{def:b3:rings:ideal}{ideais} $I_1 \subseteq I_2
\subseteq \cdots$ é estacionária.
\end{lemma}

\begin{proof}
$I = \bigcup_n I_n$ é um \omterm{def:b3:rings:ideal}{ideal} (a reunião é crescente), logo $I =
(a)$; o elemento
$a$ pertence a algum $I_N$ e, então, $I = (a)
\subseteq I_N \subseteq I_n \subseteq I$ para $n \geq N$.
\end{proof}

\begin{lemma}[Bézout; lema de Euclides]\label{lem:b3:rings:bezout}
Sejam $A$ um \omterm{def:b3:rings:pidufd}{DIP} e $a, b \in A$. Então $(a) + (b) = (d)$ para algum $d$, um
\emph{máximo divisor comum}\index{máximo divisor comum}: $d \mid a$, $d \mid b$,
e todo divisor comum de $a,
b$ divide $d$; além disso, $d = au + bv$
para certos $u, v$ (Bézout). Consequentemente, todo \omterm{def:b3:rings:divisibility}{elemento
irredutível} de um \omterm{def:b3:rings:pidufd}{DIP} é primo.
\end{lemma}

\begin{proof}
$(a) + (b)$ é um \omterm{def:b3:rings:ideal}{ideal}, logo é $(d)$; $a, b \in (d)$ dá $d \mid
a, b$; e $d = au + bv \in (a) + (b)$. Um divisor
comum $c$ de $a, b$ divide $au + bv = d$.

Euclides: seja $p$ \omterm{def:b3:rings:divisibility}{irredutível}, $p \mid ab$, $p \nmid a$. Um mdc $d$ de
$p$ e $a$ divide $p$, logo $d$ é inversível ou associado
a $p$ (irredutibilidade); o caso associado está excluído por $p \nmid a$.
Assim $1 = pu + av$, donde $b = pub + abv$, e $p$ divide ambos os termos: $p \mid b$.
\end{proof}

\begin{theorem}\label{thm:b3:rings:pidufd}
Principal $\Rightarrow$ fatorial.
\end{theorem}

\begin{proof}
\emph{Existência.} Suponha que algum elemento $a$ não nulo e não
inversível não admita fatoração em irredutíveis. Então $a$ não é
\omterm{def:b3:rings:divisibility}{irredutível}: $a =
a_1b_1$ com ambos os fatores não inversíveis; ao menos um
deles, digamos $a_1$, também não admite fatoração (um produto de
dois elementos fatoráveis é fatorável). Iterando, obtemos $a = a_0, a_1, a_2,
\dots$, cada
um divisor próprio do anterior e sem fatoração, de modo que $(a_0) \subsetneq (a_1) \subsetneq (a_2) \subsetneq \cdots$ ---
as inclusões são estritas porque $a_n = a_{n+1}c$ com $c$ não inversível
significa que $(a_n) = (a_{n+1})$ forçaria $c \in A^\times$ (cancele, no domínio). Isso
contradiz o~\cref{lem:b3:rings:acc}.

\emph{Unicidade.} Seja $p_1 \cdots p_r = q_1 \cdots q_s$ com todos os fatores irredutíveis;
mostremos $r \leq s$ por indução sobre $r$. O \omterm{def:b3:rings:divisibility}{elemento primo}
(\cref{lem:b3:rings:bezout}) $p_1$ divide o membro direito, logo divide algum
$q_j$; renumere para $j = 1$. Como $q_1$ é \omterm{def:b3:rings:divisibility}{irredutível} e
$p_1$ não é inversível, $q_1 = u p_1$ com $u \in A^\times$: $p_1,
q_1$ são \omterm{def:b3:rings:divisibility}{associados}.
Cancele $p_1$: $p_2 \cdots p_r = (u q_2)
q_3\cdots q_s$ e conclua por indução ($r = 1$ força
$s = 1$: uma unidade vezes irredutíveis não pode dar $1$).
\end{proof}

\begin{remark}\label{rem:b3:rings:ufdgcd}
Em um \omterm{def:b3:rings:pidufd}{DFU} existem mdc (tome os expoentes mínimos nas fatorações) e vale
o lema de Euclides --- \omterm{def:b3:rings:divisibility}{irredutível} $=$ primo (\cref{exo:b3:rings:2}) ---,
mas Bézout pode falhar: em $\Z[X]$, $\gcd(2, X) = 1$ e no entanto $1 \neq 2U + XV$
(avalie em $X = 0$: $1 = 2U(0)$, impossível). As identidades de Bézout
são propriedade exclusiva dos \omterm{def:b3:rings:pidufd}{DIP}.
\end{remark}

\begin{example}[Um anel sem fatoração única]
\label{ex:b3:rings:nonufd}
Nenhuma das implicações \omterm{def:b3:rings:pidufd}{euclidiano} $\Rightarrow$ \omterm{def:b3:rings:pidufd}{DIP} $\Rightarrow$ \omterm{def:b3:rings:pidufd}{DFU} é uma
equivalência, e vale a pena ver ao menos uma vez, em detalhe completo,
como a última falha. Em
\[
A = \Z[\iu\sqrt5] = \{a + \iu b\sqrt5 : a, b \in \Z\},
\qquad N(a + \iu b\sqrt5) = a^2 + 5b^2,
\]
a norma é multiplicativa e $N(z) = 1$ se, e somente se, $z \in
A^\times = \{\pm1\}$. Considere
\[
6 = 2 \cdot 3 = (1 + \iu\sqrt5)(1 - \iu\sqrt5).
\]
Os quatro fatores são irredutíveis: suas normas são $4, 9, 6,
6$, e uma
fatoração própria $z = z_1z_2$ forçaria $N(z_1) \in \{2, 3\}$ --- mas $a^2 + 5b^2$ nunca vale $2$
nem $3$ ($b = 0$ deixa os não quadrados $2, 3$; $\abs b \geq
1$ dá
$\geq 5$). E, no entanto, $2$ não é associado a nenhum dos
$1 \pm
\iu\sqrt5$ (normas $4 \neq 6$): duas fatorações genuinamente distintas de
$6$ em irredutíveis. De maneira equivalente, aqui \omterm{def:b3:rings:divisibility}{irredutível}
$\neq$ primo: $2$ divide o produto $(1 +
\iu\sqrt5)(1 - \iu\sqrt5) = 6$ mas não divide nenhum
dos fatores (normas, de novo). O reparo desse fracasso pela via dos
\omterm{def:b3:rings:ideal}{ideais} --- fatorar \emph{\omterm{def:b3:rings:ideal}{ideais}} em vez de elementos --- é a certidão
de nascimento da teoria algébrica dos números; em nosso nível, o
exemplo calibra o quanto são especiais os anéis \omterm{def:b3:rings:pidufd}{euclidianos} $\Z$,
$K[X]$ e $\Z[\iu]$ deste capítulo.
\end{example}

\begin{method}\label{met:b3:rings:quotient}
Para identificar um \omterm{def:b3:rings:ideal}{anel quociente} $A/I$, procure um morfismo
sobrejetor $f \colon A \to B$ de núcleo $I$ e invoque o~\cref{thm:b3:rings:firstiso}; quando
$A = C[X]$ é um anel de polinômios, $f$ costuma ser uma avaliação. Assim
$\Z[X]/(X^2+1) \cong
\Z[\iu]$ (avalie em $\iu$), $K[X,Y]/(Y - X^2) \cong K[X]$ (avalie $Y$ em $X^2$),
$\R[X]/(X^2+1) \cong \C$. Para mostrar que $I$ é primo ou \omterm{def:b3:rings:primemaximal}{maximal}, mostre que o
quociente é um domínio ou um corpo (\cref{prop:b3:rings:primemaximal}).
\end{method}

\section{Polinômios sobre um DFU: Gauss e Eisenstein}

Ao longo de toda esta seção, $A$ é um \omterm{def:b3:rings:pidufd}{DFU} com corpo de frações
$K$ (construído como o corpo dos quocientes formais $a/b$,
$b \neq 0$, exatamente como $\Q$ a partir de $\Z$; o volume do
segundo ano fez essa construção para $\Q$, e ela se transporta
literalmente). Nosso objetivo: a fatorialidade passa de $A$ a
$A[X]$, e a irredutibilidade sobre $A$ é essencialmente a
irredutibilidade sobre o corpo maior $K$.

\begin{definition}\label{def:b3:rings:content}
O \emph{conteúdo}\index{conteúdo de um polinômio} $c(P)$ de um
$P \in A[X]$ não nulo é um mdc de seus coeficientes (definido a menos de uma
unidade); $P$ é \emph{primitivo} se $c(P) \in A^\times$. Todo $P
\in A[X]$ se escreve
$P = c(P)\,P_1$ com $P_1$ primitivo, e todo $P \in K[X]\setminus\{0\}$ se escreve $P = \lambda P_1$ com $\lambda
\in K^\times$ e
$P_1 \in A[X]$ primitivo (elimine os denominadores e ponha o conteúdo em
evidência).
\end{definition}

\begin{lemma}[Gauss]\label{lem:b3:rings:gauss}
O produto de dois polinômios primitivos de $A[X]$ é primitivo;
consequentemente, $c(PQ) = c(P)c(Q)$ a menos de unidades.\index{lema de Gauss}
\end{lemma}

\begin{proof}
Sejam $P, Q$ primitivos e suponha que algum \omterm{def:b3:rings:divisibility}{irredutível} (= primo,
por ser \omterm{def:b3:rings:pidufd}{DFU}) $p$ divida todos os coeficientes de $PQ$. Reduza
módulo $p$: em $(A/(p))[X]$, $\bar P \bar Q = 0$. Mas $A/(p)$ é um domínio ($(p)$
é primo), logo $(A/(p))[X]$ é um domínio (os coeficientes líderes se
multiplicam), o que força $\bar P = 0$ ou $\bar Q = 0$: $p$ divide todos os
coeficientes de $P$ ou todos os de $Q$, contradizendo a
primitividade. Quanto à consequência, escreva $P = c(P)P_1$, $Q = c(Q)Q_1$: $PQ =
c(P)c(Q) P_1Q_1$ com
$P_1Q_1$ primitivo.
\end{proof}

\begin{theorem}\label{thm:b3:rings:gaussufd}
Sejam $A$ um \omterm{def:b3:rings:pidufd}{DFU} e $K$ seu corpo de frações.
\begin{enumerate}
  \item Um $P \in A[X]$ primitivo de grau $\geq 1$ é \omterm{def:b3:rings:divisibility}{irredutível} em
        $A[X]$ se, e somente se, é \omterm{def:b3:rings:divisibility}{irredutível} em $K[X]$.
  \item $A[X]$ é um \omterm{def:b3:rings:pidufd}{DFU}; seus irredutíveis são os irredutíveis de
        $A$ e os polinômios primitivos irredutíveis sobre $K$.
        Em particular, $\Z[X]$ e, por indução, $K[X_1, \dots, X_n]$ e $\Z[X_1, \dots, X_n]$ são
        \omterm{def:b3:rings:pidufd}{DFU}.
\end{enumerate}
\end{theorem}

\begin{proof}
(1) ($\Leftarrow$) Se $P = QR$ em $A[X]$ com $Q, R$ não
inversíveis, então nenhum dos fatores é constante (um fator constante
de um polinômio primitivo é uma unidade), de modo que a fatoração é
própria em $K[X]$. ($\Rightarrow$) Suponha $P = QR$ com $Q, R \in K[X]$ de graus
$\geq 1$. Escreva $Q = \lambda Q_1$, $R = \mu R_1$ com $Q_1, R_1
\in A[X]$ primitivos: $P = \lambda\mu\, Q_1R_1$, e
$Q_1R_1$ é primitivo por Gauss. Tomando \omterm{def:b3:rings:content}{conteúdos}, $\lambda\mu \in A^\times$ (ambos os
lados têm \omterm{def:b3:rings:content}{conteúdo} unidade; formalmente, $\lambda\mu = c(P) \in
A^\times$ a menos de unidades e,
em particular, $\lambda \mu \in A$): $P
= (\lambda\mu Q_1) R_1$ é uma fatoração própria em $A[X]$.

(2) Existência: dado $P \neq 0$ não inversível, fatore $P = c(P)P_1$, fatore
$c(P)$ em irredutíveis de $A$ e fatore $P_1$ no \omterm{def:b3:rings:pidufd}{DFU}
$K[X]$ como $\prod Q_i$ com $Q_i \in K[X]$ irredutíveis; escrevendo $Q_i = \lambda_i R_i$ com
$R_i \in A[X]$ primitivos (portanto irredutíveis sobre $K$ e, por (1), em
$A[X]$), o produto $\prod \lambda_i$ é uma unidade de $A$ como antes, e
$P_1 = u\prod
R_i$. Unicidade: compare a parte constante e a parte polinomial de uma
fatoração; as constantes multiplicam-se em $c(P)$ (Gauss), única
pela fatorialidade de $A$; as partes polinomiais dão duas
fatorações em $K[X]$ de um mesmo polinômio, de modo que coincidem a
menos de constantes de $K^\times$ (fatorialidade de $K[X]$,
\cref{thm:b3:rings:pidufd}), e polinômios primitivos \omterm{def:b3:rings:divisibility}{associados} em $K[X]$ são
\omterm{def:b3:rings:divisibility}{associados} em $A[X]$: se $R = \lambda R'$ com $R, R'$ primitivos e $\lambda \in K^\times$,
tomar \omterm{def:b3:rings:content}{conteúdos} força $\lambda \in A^\times$.
\end{proof}

\begin{theorem}[Critérios de irredutibilidade]\label{thm:b3:rings:criteria}
Sejam $A$ um \omterm{def:b3:rings:pidufd}{DFU}, $K$ seu corpo de frações e $P = a_nX^n + \dots
+ a_0 \in A[X]$ primitivo de
grau $n \geq 1$.
\begin{enumerate}
  \item (\emph{Redução}) Se $p \in A$ é primo, $p \nmid a_n$ e a redução
        $\bar P$ é \omterm{def:b3:rings:divisibility}{irredutível} em $(A/(p))[X]$, então $P$ é
        \omterm{def:b3:rings:divisibility}{irredutível} em $K[X]$ (e portanto em $A[X]$).
  \item (\emph{Eisenstein}\index{critério de Eisenstein}) Se algum
        primo $p$ satisfaz $p \nmid a_n$, $p \mid a_i$ para $0
        \leq i < n$, e $p^2 \nmid a_0$, então
        $P$ é \omterm{def:b3:rings:divisibility}{irredutível} em $K[X]$ (e portanto em $A[X]$).
\end{enumerate}
\end{theorem}

\begin{proof}
Pelo~\cref{thm:b3:rings:gaussufd}(1), uma fatoração própria sobre $K$ fornece
$P = QR$ com $Q, R \in A[X]$, $\deg Q, \deg R \geq 1$ (as constantes estão excluídas: seriam
unidades ou estragariam a primitividade).

(1) Reduza módulo $p$: $\bar P = \bar Q\bar R$ em $(A/(p))[X]$. Como $p \nmid a_n$ e o $\deg$ só
pode diminuir sob redução, $\deg \bar
Q = \deg Q \geq 1$ e $\deg\bar R = \deg R \geq 1$ (seus coeficientes líderes
multiplicam-se em $\bar a_n \neq 0$, logo nenhum deles cai): $\bar P$ se fatora
propriamente --- contradição.

(2) Reduza módulo $p$: $\bar Q \bar R = \bar P = \bar a_n X^n$ (todos os coeficientes inferiores
morrem). No domínio $(A/(p))[X]$, as fatorações de $cX^n$ ($c \ne 0$) são em
constantes e potências puras $c'X^k$: com efeito, se $\bar Q\bar R = \bar a_nX^n$ e,
digamos, $\bar Q$ tivesse um coeficiente não nulo em grau $< \deg\bar Q$,
tome os termos não nulos de menor grau: $\operatorname{val}(\bar Q\bar R) =
\operatorname{val}\bar Q + \operatorname{val}\bar R$ (domínio), o que deve
coincidir com $n = \deg\bar Q + \deg\bar R$, forçando $\operatorname{val} = \deg$ para ambos: os dois são monômios.
Como acima, os graus não caem, de modo que $Q$ e $R$ têm termos
constantes $Q(0), R(0)$ divisíveis por $p$ --- os dois, já que ambas as
reduções são monômios de grau $\geq 1$. Então $p^2 \mid
Q(0)R(0) = a_0$: contradição.
\end{proof}

\begin{example}\label{ex:b3:rings:eisenstein}
$X^n - p$ é \omterm{def:b3:rings:divisibility}{irredutível} sobre $\Q$ para todo primo $p$ e todo $n
\geq 1$
(\omterm{thm:b3:rings:criteria}{Eisenstein} em $p$): há polinômios irredutíveis de todos os graus
sobre $\Q$ --- em contraste flagrante com $\C$ (grau $1$,
d'Alembert--Gauss, demonstrado no~\cref{ch:b3:holomorphic}) e com $\R$ (graus
$1, 2$). O truque da \emph{translação} amplia o alcance de
\omterm{thm:b3:rings:criteria}{Eisenstein}: o \emph{polinômio ciclotômico}\index{polinômio ciclotômico}
de índice $p$, $\Phi_p = X^{p-1} +
\dots + X + 1 = \frac{X^p - 1}{X - 1}$, satisfaz
\[
\Phi_p(X + 1) = \frac{(X+1)^p - 1}{X}
= X^{p-1} + \binom{p}{1}X^{p-2} + \dots + \binom{p}{p-1},
\]
que é de \omterm{thm:b3:rings:criteria}{Eisenstein} em $p$ ($p \mid \binom pk$ para $0 < k < p$, e $\binom{p}{p-1} = p \not\equiv 0 \bmod p^2$): logo $\Phi_p(X+1)$
e, com ele, $\Phi_p$ são irredutíveis sobre $\Q$. Esse é o
coração algébrico da história do $17$-ágono contada no~\cref{ch:b3:galois}.
\end{example}

\begin{method}\label{met:b3:rings:irreducibility}
Para demonstrar que $P \in \Z[X]$ é \omterm{def:b3:rings:divisibility}{irredutível} sobre $\Q$: (i) torne
$P$ primitivo; (ii) tente \omterm{thm:b3:rings:criteria}{Eisenstein}, em $P(X)$ e nas
translações $P(X \pm
1)$; (iii) tente a redução módulo primos pequenos que não
dividam o coeficiente líder --- a irredutibilidade módulo \emph{um}
único $p$ já basta, e sobre $\mathbb F_p$ a irredutibilidade é uma
verificação finita (a ausência de raízes exclui fatores de grau
$1$; depois, teste os finitos fatores de cada grau $\leq \deg P/2$); (iv) se
tudo o mais falhar, coeficientes indeterminados. Cuidado: ser redutível
módulo todo $p$ \emph{não} implica ser redutível sobre $\Q$
(\cref{exo:b3:rings:11}).
\end{method}

\section{Anéis noetherianos}

\begin{definition}\label{def:b3:rings:noetherian}
Um anel $A$ é \emph{noetheriano}\index{anel noetheriano} se todo
\omterm{def:b3:rings:ideal}{ideal} de $A$ é finitamente gerado.
\end{definition}

\begin{proposition}\label{prop:b3:rings:noethacc}
$A$ é \omterm{def:b3:rings:noetherian}{noetheriano} se, e somente se, toda sequência crescente de
\omterm{def:b3:rings:ideal}{ideais} é estacionária (\emph{condição de cadeia ascendente}), e se, e
somente se, toda família não vazia de \omterm{def:b3:rings:ideal}{ideais} tem elemento \omterm{def:b3:rings:primemaximal}{maximal} (para
a inclusão).
\end{proposition}

\begin{proof}
\emph{(Finitamente gerado $\Rightarrow$ CCA)}: para uma cadeia $I_1 \subseteq I_2
\subseteq \cdots$, a reunião
$I$ é um \omterm{def:b3:rings:ideal}{ideal}, gerado por $x_1,
\dots, x_r$; todos os $x_i$ pertencem a
algum $I_N$, logo $I = I_N = I_n$ para $n \geq N$. \emph{(CCA $\Rightarrow$ elementos
maximais)}: se uma família não vazia $\mathcal F$ não tivesse elemento
\omterm{def:b3:rings:primemaximal}{maximal}, escolha $I_1 \in
\mathcal F$ e, indutivamente, $I_{n+1} \supsetneq I_n$ em $\mathcal F$ (é possível,
pois $I_n$ não é \omterm{def:b3:rings:primemaximal}{maximal}): uma cadeia infinita estritamente
crescente. (Isso usa o axioma das escolhas dependentes, uma forma fraca
da escolha com a qual não nos preocupamos.) \emph{(Elementos maximais
$\Rightarrow$ finitamente gerado)}: dado um \omterm{def:b3:rings:ideal}{ideal} $I$, a família dos \omterm{def:b3:rings:ideal}{ideais}
finitamente gerados contidos em $I$ é não vazia ($(0)$); um
elemento \omterm{def:b3:rings:primemaximal}{maximal} $J = (x_1, \dots, x_r)$ tem de ser igual a $I$: caso contrário,
acrescentar $x \in I \setminus J$ aos geradores produziria um membro estritamente maior
da família.
\end{proof}

\begin{theorem}[Teorema da base de Hilbert]\label{thm:b3:rings:hilbert}
Se $A$ é \omterm{def:b3:rings:noetherian}{noetheriano}, $A[X]$ também é. Logo $A[X_1, \dots,
X_n]$ também são, e
todo quociente deles.\index{teorema da base de Hilbert}
\end{theorem}

\begin{proof}
Seja $I$ um \omterm{def:b3:rings:ideal}{ideal} de $A[X]$ e suponha que $I$ não seja
finitamente gerado. Construa uma sequência: $f_1 \in I \setminus \{0\}$ de grau mínimo e,
indutivamente, $f_{k+1} \in I \setminus (f_1,
\dots, f_k)$ de grau mínimo (o conjunto é não vazio, por
hipótese). Os graus $d_k = \deg f_k$ são não decrescentes (pela minimalidade de
cada escolha: $f_{k+1}$ estava disponível na etapa $k+1$\dots\ mais
precisamente, $f_{k+1} \notin (f_1,\dots,f_k) \supseteq
(f_1, \dots, f_{k-1})$, de modo que $f_{k+1}$ concorreu na etapa $k$ e
perdeu ou empatou: $d_{k+1} \geq d_k$). Seja $a_k \in A$ o coeficiente líder de $f_k$.
A cadeia de \omterm{def:b3:rings:ideal}{ideais} $(a_1) \subseteq (a_1,
a_2) \subseteq \cdots$ estabiliza: $a_{n+1} \in (a_1, \dots,
a_n)$ para algum $n$, digamos
$a_{n+1} = \sum_{k\leq n} u_k a_k$. Considere
\[
g = f_{n+1} - \sum_{k=1}^{n} u_k X^{\,d_{n+1} - d_k} f_k .
\]
Então $g \in I \setminus (f_1, \dots, f_n)$ (a soma pertence ao \omterm{def:b3:rings:ideal}{ideal}, $f_{n+1}$ não), e no entanto o
coeficiente de grau $d_{n+1}$ se cancela: $\deg g < d_{n+1}$, contradizendo a
minimalidade de $d_{n+1} = \deg f_{n+1}$.

Iterando, $A[X_1, \dots, X_n] = (A[X_1, \dots, X_{n-1}])[X_n]$ é \omterm{def:b3:rings:noetherian}{noetheriano}; um quociente $A/I$ é \omterm{def:b3:rings:noetherian}{noetheriano}
porque seus \omterm{def:b3:rings:ideal}{ideais} $J/I$ se levantam a \omterm{def:b3:rings:ideal}{ideais} de $A$
(correspondência), onde uma quantidade finita de geradores se projeta
sobre geradores.
\end{proof}

\begin{remark}\label{rem:b3:rings:noethmeaning}
A noetherianidade é o axioma de finitude da geometria algébrica:
qualquer sistema de equações polinomiais em $n$ variáveis, por mais
infinito que seja, é equivalente a um número finito delas --- seu
conjunto de soluções é recortado por finitos polinômios. Os \omterm{def:b3:rings:pidufd}{DIP} são
\omterm{def:b3:rings:noetherian}{noetherianos} (trivialmente); $\Z[X_1, X_2, \dots]$ em uma infinidade de variáveis não é
($(X_1) \subsetneq (X_1, X_2) \subsetneq \cdots$). Anéis não \omterm{def:b3:rings:noetherian}{noetherianos} também surgem naturalmente em análise:
as funções contínuas sobre $\intcc01$ formam um deles (\cref{exo:b3:rings:10}).
\end{remark}

\section{Exercícios}

\begin{exercise}[$\star$]\label{exo:b3:rings:1}
Identifique os quocientes: (a) $\Z[X]/(X^2 + 1) \cong \Z[\iu]$; (b) $\R[X]/(X^2+1) \cong \C$; (c) $\mathbb F_2[X]/(X^2 + X + 1)$ é um corpo
com $4$ elementos --- escreva sua tabela de multiplicação.
\end{exercise}

\begin{exercise}[$\star$]\label{exo:b3:rings:2}
(a) Mostre que, em um \omterm{def:b3:rings:pidufd}{DFU}, todo \omterm{def:b3:rings:divisibility}{elemento irredutível} é primo.
(b) Mostre que um domínio de integridade finito é um corpo.
(c) Deduza que, em um anel finito, todo \omterm{def:b3:rings:primemaximal}{ideal primo} é \omterm{def:b3:rings:primemaximal}{maximal}.
\end{exercise}

\begin{exercise}[$\star$]\label{exo:b3:rings:3}
Em $\Z[\iu\sqrt5]$: verifique que $3$, $1 + \iu\sqrt5$ e $1 -
\iu\sqrt5$ são irredutíveis, que
$9 = 3\cdot 3 = (2 +
\iu\sqrt5)(2 - \iu\sqrt5)$, e conclua de novo (depois da~\cref{prop:b3:rings:primeirred}) que $\Z[\iu\sqrt5]$ não é um
\omterm{def:b3:rings:pidufd}{DFU}. Onde exatamente a unicidade falha?
\end{exercise}

\begin{exercise}[$\star\star$]\label{exo:b3:rings:4}
(a) Mostre que $\Z[\iu]$ é \omterm{def:b3:rings:pidufd}{euclidiano} para a norma $N(x + \iu y) =
x^2 + y^2$: dados $a, b \neq 0$,
escolha $q \in \Z[\iu]$ mais próximo de $a/b \in \C$.
(b) Determine $\Z[\iu]^\times$.
(c) Mesmas questões para $\Z[\iu\sqrt2]$ e $N(x + \iu y\sqrt2) =
x^2 + 2y^2$. Por que o mesmo argumento falha
para $\Z[\iu\sqrt5]$?
\end{exercise}

\begin{exercise}[$\star\star$]\label{exo:b3:rings:5}
Seja $A$ um anel. (a) Mostre que, se $x$ é nilpotente ($x^n = 0$
para algum $n$), então $1 + x \in A^\times$. (b) Mostre que, se $A$ é um
domínio, $A[X]^\times = A^\times$; dê um contraexemplo sobre $\Z/4\Z$. (c) Mostre que
um domínio não tem idempotentes ($e^2 = e$) além de $0, 1$, nem
nilpotentes além de $0$.
\end{exercise}

\begin{exercise}[$\star\star$]\label{exo:b3:rings:6}
Em $A = K[X, Y]$: (a) mostre que o \omterm{def:b3:rings:ideal}{ideal} $(X, Y)$ é \omterm{def:b3:rings:primemaximal}{maximal} mas não é
principal --- de modo que $K[X,Y]$ é um \omterm{def:b3:rings:pidufd}{DFU} (\cref{thm:b3:rings:gaussufd}) que
não é um \omterm{def:b3:rings:pidufd}{DIP}; (b) identifique $K[X, Y]/(Y - X^2)$ e $K[X,Y]/(XY - 1)$ como subanéis de funções
racionais; (c) $(Y - X^2)$ é primo? é \omterm{def:b3:rings:primemaximal}{maximal}?
\end{exercise}

\begin{exercise}[$\star\star$]\label{exo:b3:rings:7}
\omterm{def:b3:rings:divisibility}{Irredutível} ou não sobre $\Q$: $X^5 - 12X^3 + 36X - 12$;\quad $X^4 + X + 1$ \emph{(reduza
módulo $2$)};\quad $X^4 + 4$;\quad $\Phi_8 = X^4 + 1$ \emph{(translade por
$1$)};\quad $X^3 - X - 1$.
\end{exercise}

\begin{exercise}[$\star\star$]\label{exo:b3:rings:8}
(a) A partir do~\cref{thm:b3:rings:crt}, demonstre que a função de Euler é
multiplicativa em argumentos coprimos e que $\varphi(p^k) =
p^{k-1}(p-1)$; recupere $\varphi(n) = n\prod_{p \mid n}(1 -
\frac1p)$.
(b) Resolva: $x \equiv 2 \pmod 7$, $x \equiv 5 \pmod{11}$, $x
\equiv 1 \pmod{13}$, exibindo os idempotentes $e_k$ da
demonstração do~\cref{thm:b3:rings:crt}.
\end{exercise}

\begin{exercise}[$\star\star\star$]\label{exo:b3:rings:9}
(O nilradical) Seja $\operatorname{Nil}(A)$ o conjunto dos elementos nilpotentes. (a)
Mostre que $\operatorname{Nil}(A)$ é um \omterm{def:b3:rings:ideal}{ideal} contido em todo \omterm{def:b3:rings:primemaximal}{ideal primo}. (b)
Reciprocamente, seja $a$ não nilpotente; usando o lema de Zorn
sobre os \omterm{def:b3:rings:ideal}{ideais} que evitam $S =
\{a^n : n \in \N\}$, produza um \omterm{def:b3:rings:primemaximal}{ideal primo} que não
contenha $a$. Conclua:
\[
\operatorname{Nil}(A) = \bigcap_{\mathfrak p \text{ primo}}
\mathfrak p .
\]
\end{exercise}

\begin{exercise}[$\star\star\star$]\label{exo:b3:rings:10}
(a) Sejam $A$ \omterm{def:b3:rings:noetherian}{noetheriano} e $f \colon A \to A$ um morfismo de anéis sobrejetor.
Mostre que $f$ é injetor. \emph{(Considere $\ker
f \subseteq \ker f^2 \subseteq \cdots$.)}
(b) Mostre que o anel $\mathcal C(\intcc01, \R)$ das funções contínuas não é \omterm{def:b3:rings:noetherian}{noetheriano}.
\emph{(Considere $I_n = \{f : f = 0
\text{ em } \intcc0{1/n}\}$.)}
(c) Mostre que, em um \emph{domínio} \omterm{def:b3:rings:noetherian}{noetheriano}, todo elemento não
nulo e não inversível é um produto (finito) de irredutíveis --- de modo
que a não fatorialidade de $\Z[\iu\sqrt 5]$ é uma falha apenas de
\emph{unicidade}.
\end{exercise}

\begin{exercise}[$\star\star\star$]\label{exo:b3:rings:11}
Seja $P = X^4 + 1$.
(a) Mostre que $P$ é \omterm{def:b3:rings:divisibility}{irredutível} sobre $\Q$ (\cref{exo:b3:rings:7}).
(b) Mostre que $P$ é redutível módulo \emph{todo} primo $p$:
trate $p = 2$; depois, para $p$ ímpar, mostre que $8 \mid p^2 - 1$ e
admita por ora (demonstrado no~\cref{ch:b3:galois}) que o grupo
multiplicativo do corpo com $p^2$ elementos é cíclico, para concluir
que $P$ se decompõe em dois fatores quadráticos módulo $p$;
explicite-os quando um dos números $-1$, $2$, $-2$ é um
quadrado módulo $p$, e mostre que um deles sempre é.
\end{exercise}

\begin{exercise}[$\star\star$]\label{exo:b3:rings:12}
(Idempotentes decompõem anéis) Um elemento $e$ de um anel comutativo
$A$ é \emph{idempotente} se $e^2 = e$.
(a) Mostre que, se $e$ é idempotente, $1 - e$ também é, e que a
aplicação $x \mapsto (ex, (1-e)x)$ é um isomorfismo de anéis $A
\cong Ae \times A(1-e)$, onde $Ae$ é um anel
de unidade $e$.
(b) Encontre todos os idempotentes de um domínio e os de $\Z/12\Z$; exiba o
isomorfismo $\Z/12\Z \cong \Z/4\Z \times \Z/3\Z$ nomeando seus dois idempotentes não triviais.
(c) Mostre que a decomposição de $\Z/n\Z$ pelo teorema chinês dos
restos (\cref{ex:b3:rings:crtz}) corresponde exatamente aos idempotentes $e_i \equiv 1 \bmod p_i^{a_i}$,
$e_i \equiv 0$ módulo as demais potências de primos: os anéis se decompõem ao
longo de seus idempotentes como os espaços se decompõem ao longo de
projeções.
\end{exercise}

\section{Problema: o teorema dos dois quadrados de Fermat}

\begin{problem}[{Problema de fim de semana --- somas de dois quadrados,
via $\Z[\iu]$}]\label{pb:b3:rings:1}
Quais inteiros são somas de dois quadrados? A resposta de Fermat (1640)
é uma das joias da aritmética; os inteiros de Gauss transformam sua
demonstração em teoria de anéis. Ao longo de todo o problema, $N(x + \iu y) = x^2 + y^2$
designa a norma, $\Z[\iu]$ é \omterm{def:b3:rings:pidufd}{euclidiano} (\cref{exo:b3:rings:4}), logo um \omterm{def:b3:rings:pidufd}{DIP} e um
\omterm{def:b3:rings:pidufd}{DFU}, e \emph{primo de Gauss} significa \omterm{def:b3:rings:divisibility}{elemento primo} (= \omterm{def:b3:rings:divisibility}{irredutível})
de $\Z[\iu]$.

\medskip
\noindent\textbf{Parte I --- Normas e primos de Gauss.}
\begin{enumerate}
  \item Verifique $N(zw) = N(z)N(w)$, deduza de novo $\Z[\iu]^\times =
        \{\pm1, \pm\iu\}$ e demonstre a
        \emph{identidade de Brahmagupta}: um produto de duas somas de
        dois quadrados é uma soma de dois quadrados.
  \item Mostre que, se $N(z)$ é um número primo, então $z$ é um
        primo de Gauss.
  \item Mostre que todo primo de Gauss $\pi$ divide exatamente um
        número primo $p$ \emph{(considere $N(\pi) = \pi\bar\pi$)} e que, então,
        $N(\pi) \in \{p, p^2\}$.
  \item Deduza a dicotomia: para cada primo $p$, ou $p$
        permanece primo em $\Z[\iu]$ (e não existe primo de Gauss de norma
        $p$), ou $p = \pi\bar\pi$ com $\pi$ um primo de Gauss de norma
        $p$ --- e então $p = a^2 + b^2$.
\end{enumerate}

\medskip
\noindent\textbf{Parte II --- Teorema de Wilson e $-1$ módulo
$p$.}
\begin{enumerate}[resume]
  \item Demonstre o \emph{teorema de Wilson}: para $p$ primo,
        $(p-1)!
        \equiv -1 \pmod p$. \emph{(Emparelhe cada resto com seu inverso; quais são
        os que se emparelham consigo mesmos?)}
  \item Sejam $p$ um primo ímpar e $m = \frac{p-1}2$. Mostre que $(m!)^2 \equiv (-1)^{m+1} \pmod p$
        \emph{(em $(p-1)!$, substitua cada fator $k > m$ por
        $-(p - k)$)}.
  \item Conclua: $-1$ é um quadrado módulo $p$ se, e somente se,
        $p = 2$ ou $p
        \equiv 1 \pmod 4$. \emph{(Para o ``somente se'': se $x^2
        \equiv -1$,
        qual é a ordem de $x$ em $(\Z/p\Z)^\times$, e o que diz Lagrange?)}
\end{enumerate}

\medskip
\noindent\textbf{Parte III --- A lei de decomposição.}
\begin{enumerate}[resume]
  \item Sejam $p \equiv 1 \pmod 4$ e $x$ com $p \mid x^2 + 1 =
        (x + \iu)(x - \iu)$. Mostre que $p$ \emph{não}
        é um primo de Gauss e conclua, com a Parte I: $p = a^2 + b^2$.
  \item Seja $p \equiv 3 \pmod 4$. Mostre diretamente que $p$ não é uma soma de
        dois quadrados \emph{(quadrados módulo $4$)} e deduza que
        $p$ permanece um primo de Gauss.
  \item Resolva o caso $p = 2$: exiba a fatoração $2 =
        -\iu(1+\iu)^2$ e verifique
        que $1 + \iu$ é um primo de Gauss. ($2$ é o único primo
        \emph{ramificado}: divisível pelo quadrado de um primo de Gauss
        a menos de uma unidade.)
  \item Monte a classificação dos primos de Gauss, a menos de unidades:
        $1 + \iu$; os inteiros $p \equiv 3 \pmod 4$; os pares conjugados $\pi, \bar\pi$ de norma
        $p \equiv 1 \pmod
        4$. Verifique-a em $5 = (2+\iu)(2-\iu)$ e em $3$.
\end{enumerate}

\medskip
\noindent\textbf{Parte IV --- O teorema dos dois quadrados.}
\begin{enumerate}[resume]
  \item Demonstre a metade direta: se, na fatoração $n =
        \prod p_i^{\alpha_i}$, todo primo
        $\equiv 3 \pmod 4$ aparece com expoente par, então $n$ é uma soma de dois
        quadrados. \emph{(Brahmagupta + Partes II--III.)}
  \item Demonstre a recíproca: se $n = a^2 + b^2 = N(a + \iu b)$ e $q \equiv 3 \pmod 4$ divide $n$, mostre
        que $q$, primo de Gauss, divide $a + \iu b$ ou $a - \iu b$, que ele de
        fato divide tanto $a$ quanto $b$, e conclua por indução
        sobre $n$ que o expoente de $q$ em $n$ é par.
  \item Enuncie o teorema final. Quais dentre $2025$, $2026$,
        $2027$ são somas de dois quadrados? \emph{($2025 = 81 \cdot 25$; $2026 = 2 \cdot 1013$,
        com $1013$ primo; $2027$ primo.)}
  \item (Epílogo) Mostre que um primo $p \equiv 1 \pmod 4$ é uma soma de dois
        quadrados de maneira essencialmente única: se $p =
        a^2 + b^2 = c^2 + d^2$ (inteiros
        positivos), então $\{a, b\}
        = \{c, d\}$. \emph{(Unicidade da fatoração em
        $\Z[\iu]$.)}
\end{enumerate}

\medskip
\noindent\textbf{Parte V --- Contando representações: a fórmula de
Jacobi e a série de Leibniz.} Escreva $r_2(n) = \#\{(a, b) \in
\Z^2 : a^2 + b^2 = n\}$ (pares ordenados, sinais e
zeros incluídos) e seja $\chi$ o caractere não trivial módulo
$4$: $\chi(d) = +1$ se $d \equiv 1$, $-1$ se $d \equiv 3 \pmod4$, $0$ se $d$ é par.
\begin{enumerate}[resume]
  \item (Aquecimento, por contraste) Quais inteiros são
        \emph{diferenças} de dois quadrados? Mostre: $n = a^2 - b^2$ com $a, b \in \Z$ se,
        e somente se, $n \not\equiv 2 \pmod 4$ --- sem nenhuma teoria de anéis e sem
        estrutura comparável à que segue.
  \item Mostre que $r_2(n)$ é o número de $z \in \Z[\iu]$ com $N(z) = n$.
        Escrevendo $n = 2^{a}\prod_jp_j^{b_j}
        \prod_kq_k^{c_k}$ com $p_j \equiv 1$, $q_k \equiv 3
        \pmod4$, use a classificação da
        questão 11 e a fatoração única para mostrar: tais $z$
        existem se, e somente se, todos os $c_k$ são pares e, nesse
        caso,
        \[
        r_2(n) = 4\prod_j\,(b_j + 1) .
        \]
        \emph{(Conte: $z = u\,(1+\iu)^{a}\prod_j\pi_j^{s_j}
        \bar\pi_j^{\,b_j - s_j}\prod_kq_k^{c_k/2}$ com $u$ unidade e $0 \leq s_j \leq b_j$; por que essa
        lista é exaustiva e sem repetições?)}
  \item Mostre que $d \mapsto \chi(d)$ é completamente multiplicativa, deduza que
        $n \mapsto
        \sum_{d \mid n}\chi(d)$ é multiplicativa e calcule-a nas potências de primos: ela
        vale $1$ em $2^a$; $b + 1$ em $p^b$ ($p \equiv 1$);
        $1$ ou $0$ em $q^c$ ($q \equiv 3$), conforme $c$ seja
        par ou ímpar.
  \item Conclua o \emph{teorema de Jacobi}:
        \[
        r_2(n) = 4\sum_{d \mid n}\chi(d) =
        4\bigl(d_1(n) - d_3(n)\bigr),
        \]
        onde $d_i(n)$ conta os divisores $\equiv i \pmod 4$. Verifique em $n = 3, 5, 9, 25$ e
        liste as $16$ representações de $65$.
  \item (O círculo) Mostre que $\sum_{n \leq x}r_2(n)$ é o número de pontos do reticulado
        $\Z^2$ no disco fechado de raio $\sqrt x$, e demonstre
        \[
        \sum_{n\leq x}r_2(n) = \pi x + O(\sqrt x)
        \]
        \emph{(cada ponto do reticulado é dono de um quadrado unitário;
        compare áreas, com o erro vivendo em um anel de largura
        $O(1)$)}.
  \item (Leibniz, lido aritmeticamente) Combine as questões 19--20:
        \[
        \sum_{d \leq x}\chi(d)\Bigl\lfloor\frac
        xd\Bigr\rfloor = \frac{\pi x}4 + O(\sqrt x),
        \]
        e deduza --- removendo as partes inteiras com cuidado --- a
        série de Leibniz
        \[
        1 - \frac13 + \frac15 - \frac17 + \dots = \frac\pi4 .
        \]
        A série alternada dos inversos ímpares \emph{é} o excesso médio
        dos divisores $\equiv 1$ sobre os divisores $\equiv 3$: a análise
        calculada pela aritmética.
  \item (Quão raras são as somas de dois quadrados?) Mostre que nenhum
        inteiro $\equiv 3 \pmod 4$ é soma de dois quadrados (de duas maneiras:
        quadrados módulo $4$, ou o critério de paridade da questão
        17), de modo que ao menos um quarto de todos os inteiros fica
        de fora; e mostre que a média $\frac1x\sum_{n\leq
        x}r_2(n) \to \pi$ da questão 20 é compatível
        com o fato de os inteiros representáveis terem densidade
        $0$ --- exiba inteiros com um número anormalmente grande de
        representações (tome produtos de muitos primos $\equiv 1 \bmod 4$) para
        explicar como uma proporção que tende a zero ainda pode
        sustentar uma média positiva. (Landau demonstrou que a
        densidade verdadeira decai como $1/\sqrt{\log x}$; isso está além de nossas
        ferramentas, mas o mecanismo já se tornou visível.)
\end{enumerate}

\medskip
\noindent\textbf{Parte VI --- Complementos: representações primitivas e
Pitágoras.}
\begin{enumerate}[resume]
  \item Diga que uma representação $n = a^2 + b^2$ é \emph{primitiva} se $\gcd(a, b) = 1$.
        Mostre que $n \geq 1$ admite uma representação primitiva se, e
        somente se, $4 \nmid n$ e nenhum primo $q \equiv 3 \pmod 4$ divide $n$.
        \emph{(Para a necessidade, reutilize a descida da questão 13 e
        os quadrados módulo $4$; para a suficiência, construa
        $z$ a partir de $1 + \iu$ e dos $\pi_j$ apenas --- sem
        conjugados --- e explique por que um fator primo comum a
        $a$ e $b$ forçaria $\pi_j$ e $\bar\pi_j$, ou $(1+\iu)^2$, a
        entrarem ambos em $z$.)}
  \item (Ternos pitagóricos) Sejam $a^2 + b^2 = c^2$ com $a, b,
        c$ positivos, $\gcd(a, b) = 1$ e
        $b$ par. Mostre que $a + \iu b$ e $a - \iu b$ são coprimos em $\Z[\iu]$
        \emph{(um divisor primo de Gauss comum dividiria $2a$ e
        $2b$, e $c$ é ímpar)}, deduza da fatoração única que
        $a + \iu b = u(m + \iu n)^2$ para uma unidade $u$, e conclua a parametrização
        clássica: a menos de trocar $a$ e $b$,
        \[
        a = m^2 - n^2, \qquad b = 2mn, \qquad c = m^2 + n^2,
        \]
        com $m > n \geq 1$ coprimos de paridades opostas. Recupere $(3, 4, 5)$ e $(21, 20, 29)$
        a partir de $(m, n) =
        (2, 1)$ e de $(5, 2)$.
  \item (Verificação numérica) Tome $x = 25$. Calcule $r_2(n)$
        para $1 \leq n \leq 25$ pela fórmula de Jacobi, verifique que os valores não
        nulos ocorrem exatamente em $n = 1,
        2, 4, 5, 8, 9, 10, 13, 16, 17, 18, 20, 25$, e que
        \[
        \sum_{n \leq 25} r_2(n) = 80
        = 4\sum_{d \leq 25}\chi(d)
        \Bigl\lfloor\frac{25}d\Bigr\rfloor .
        \]
        Verifique que o disco fechado de raio $5$ contém $81$
        pontos do reticulado e compare com $\pi x \approx 78.5$: o erro está bem
        dentro do $O(\sqrt x)$ da questão 20.
\end{enumerate}
\end{problem}
